\documentclass[a4paper,10pt]{article}
\usepackage[utf8]{inputenc}
\usepackage[T2A,T1]{fontenc}
\usepackage[russian,english]{babel}
\usepackage[left=1.5cm,right=1.5cm,top=1.2cm,bottom=1.2cm]{geometry}
\usepackage{booktabs}
\usepackage{array}
\usepackage{xcolor}
\usepackage{colortbl}
\usepackage{tabularx}

\pagestyle{empty}
\renewcommand{\arraystretch}{1.25}

\definecolor{rowgrey}{HTML}{F3F3F3}
\definecolor{accent}{HTML}{2E7D32}
\definecolor{headerblue}{HTML}{1565C0}

\begin{document}
\selectlanguage{russian}

\begin{center}
{\Large\bfseries Сравнение моделей для детекции аномалий в 223-ФЗ}\\[2pt]
{\small Detection of Anomalies and Unfair Competition in Commercial Tenders Using NLP}\\[2pt]
{\small Корпус: 400 лотов, 224 manually-annotated флага.}
\end{center}

\vspace{2pt}

\begin{tabularx}{\textwidth}{>{\bfseries}lcccccX}
\toprule
\rowcolor{headerblue!90}
\multicolumn{1}{>{\bfseries\color{white}}l}{Модель} &
\multicolumn{1}{>{\color{white}}c}{Spec\textsuperscript{*}} &
\multicolumn{1}{>{\color{white}}c}{Flags\textsuperscript{**}} &
\multicolumn{1}{>{\color{white}}c}{Precision\textsuperscript{†}} &
\multicolumn{1}{>{\color{white}}c}{Cost\textsuperscript{‡}} &
\multicolumn{1}{>{\color{white}}c}{Time/lot} &
\multicolumn{1}{>{\color{white}}X}{Заметки} \\
\midrule
v1 baseline       & 5\%   & 91 & ${\sim}13\%$  & \$2.5  & 8s  & gpt-4o-mini, codebook v1 (EU/OECD) \\
\rowcolor{rowgrey}
v3+2stage         & 5\%   & 88 & ${\sim}15\%$  & \$5    & 30s & + whitelist W1--W15 + 2-stage \\
Rule detector     & 75\%  & 6  & \textbf{100\%} (6/6) & \$0    & 1s  & 5 regex-правил, deterministic \\
\rowcolor{rowgrey}
Ensemble v2$\cap$v3 & 95\% & 1 & 28.6\% (на 30 lots) & \$10 & 60s & пересечение LLM-пайплайнов \\
gpt-5.5 zero-shot & n/a    & 4 (signal) & 100\% (4/4) & \$80   & 30s & frontier, дорого \\
\rowcolor{rowgrey}
FT v1 (gpt-4.1)   & 80\%  & 5  & 60\% strict / 100\% loose & \$5  & 5s  & 122 ex, val loss 0.17 \\
\rowcolor{accent!15}
\textbf{FT v2 (gpt-4.1)} & \textbf{65\%} & \textbf{8} & \textbf{87.5\% (7/8)} & \textbf{\$5} & \textbf{5s} & \textbf{204 ex, val loss 0.03}\\
\bottomrule
\end{tabularx}

\smallskip
{\footnotesize
\textsuperscript{*}\textbf{Specificity} — \% правильно идентифицированных «чистых» лотов на 20 held-out лотах не из training.\\
\textsuperscript{**}\textbf{Flags} — общее число флагов на тех же 20 held-out лотах. Меньше = меньше шума.\\
\textsuperscript{†}\textbf{Precision} — strict (только TRUE) на held-out, manually verified для FT v1, FT v2, Rule detector. Для v1/v3+2s оценка экстраполирована с 121-flag manual review (см. \texttt{docs/results.md}).\\
\textsuperscript{‡}\textbf{Cost} — стоимость анализа 200 лотов в USD. Inference + (для FT) разовая тренировка.}

\vspace{6pt}
\noindent\textbf{Описания моделей:}

\begin{itemize}\setlength{\itemsep}{1pt}
\item \textbf{v1 baseline} --- стартовая prompt-engineered модель на gpt-4o-mini. Codebook v1 — таксономия 8 меток на основе EU Directive 2014/24/EU и OECD Bid-Rigging Indicators. Помечает почти все 223-ФЗ лоты как HIGH из-за обязательных формальных норм закона. \textit{Главный научный результат: показала, что наивный перенос EU/OECD codebook на 223-ФЗ не работает.}

\item \textbf{v3+2stage} --- prompt-engineering iteration: codebook расширен whitelist'ом W1--W15 (15 категорий стандартных норм 223-ФЗ, которые НЕ нужно флагать). Добавлен второй LLM-вызов для валидации каждого флага. Чуть лучше дискриминирует, но всё ещё over-flags на новых данных.

\item \textbf{Rule detector} --- 5 regex-правил (R1: «Не предусмотрено» в критериях; R3: прямой бан «или эквивалент»; и др.). \textbf{Manually verified: 100\% strict precision} (6/6 флагов TRUE на held-out). Дешёвый pre-filter, но узкое покрытие паттернов.

\item \textbf{Ensemble v2$\cap$v3} --- пересечение флагов из двух LLM-пайплайнов через fuzzy span matching. Самая консервативная модель: 28.6\% strict precision (manual), но recall очень низкий (7 флагов на 30 лотов).

\item \textbf{gpt-5.5 zero-shot} --- frontier модель OpenAI (release 2026-04-23). 100\% lot-level precision на signal subset (4 из 4 — TRUE), но \$80/200 лотов делает её \textbf{нерентабельной для production} (1000+ лотов/день).

\item \textbf{FT v1} --- зафайнтюненая gpt-4.1 на 122 manually-annotated examples. 80\% specificity на held-out, \textbf{manually verified: 60\% strict / 100\% loose precision} (3 TRUE + 2 PARTIAL из 5 флагов). PARTIAL флаги --- compatibility-justified brand (W7): закупка ремонта оборудования конкретного производителя НПП «ЭПРО».

\item \textbf{FT v2} --- та же база (gpt-4.1), но обучена на расширенном датасете 204 examples с балансом 46\% positive (vs v1 14\%). Источники: 121 v1-разметка + 74 manual review held-out + 78 rule-based silver labels (R1+R3) + 7 ensemble verdicts + 4 gpt-5.5 silver. Validation loss 0.03 (\textbf{5$\times$ лучше v1}). \textbf{Manual verified: 87.5\% strict precision} (7/8 флагов TRUE на held-out). Единственный FP — модель неправильно интерпретировала «*-- или эквивалент» как нарушение, тогда как это корректное обозначение в TZ. \textbf{Demonstrated transfer learning}: сгенерализовала R3 паттерн на 5 новых лотов. \textbf{Самая сильная модель в линейке}.
\end{itemize}

\vspace{4pt}
\noindent\textbf{Production cascade architecture}:
$$\text{Lot} \;\rightarrow\; \text{Rule detector (free)} \;\rightarrow\; \text{FT v2 (\$5/200)} \;\rightarrow\; \text{Human review (top flagged)}$$
$$\text{Cost: }\sim\$0.025\text{/лот}, \text{ throughput: }\sim 12\text{ лотов/мин}, \text{ scalable до 17K лотов/день на одном API key}.$$

\vspace{2pt}
\hrule
\vspace{2pt}
{\footnotesize Корпус: 400 лотов 223-ФЗ из Marker Interfax. Manual annotations: 224 верификации.
Сгенерировано: 2026-05-06.}

\end{document}
